\section*{Chapitre \ref{ch:m3:forward-curves-storage-and-convenience-yield} --- Courbes à terme, stockage et rendement d'opportunité}

\begin{solution}{exo:m3:forward-curves-storage-and-convenience-yield:1}
$80\,e^{(0.04 + 0.03 - 0.05)} = 80\,e^{0.02} = \$81.62$.
\end{solution}

\begin{solution}{exo:m3:forward-curves-storage-and-convenience-yield:2}
\omterm{def:m3:forward-curves-storage-and-convenience-yield:carry}{Portage complet}~: $80\,(e^{0.07/12} - 1) = \$0.47$ par baril.
\end{solution}

\begin{solution}{exo:m3:forward-curves-storage-and-convenience-yield:3}
Un contango plus pentu que le \omterm{def:m3:forward-curves-storage-and-convenience-yield:carry}{portage complet} est résorbé par le \omterm{def:m3:forward-curves-storage-and-convenience-yield:carry}{cash-and-carry}, que
tout détenteur de capacités de stockage peut mettre en œuvre. Arbitrer la backwardation suppose de vendre la matière maintenant et de la racheter plus tard, ce que seuls
les détenteurs de stocks peuvent le faire, et ils les gardent parce qu'ils en ont besoin~: le
\omterm{def:m3:forward-curves-storage-and-convenience-yield:storage}{rendement d'opportunité} n'a pas de borne supérieure.
\end{solution}

\begin{solution}{exo:m3:forward-curves-storage-and-convenience-yield:4}
$0.05 - 12\ln(71.20/70.00) = -15.4\%$ par an~: la courbe est en contango, d'un écart supérieur au financement, donc le
marché paie le stockage~; le \omterm{def:m3:forward-curves-storage-and-convenience-yield:storage}{rendement d'opportunité} est proche de zéro et les inventaires sont abondants.
\end{solution}

\begin{solution}{exo:m3:forward-curves-storage-and-convenience-yield:5}
$5.55 - 3.08 - 3.18 = -0.71\%$ par an.
\end{solution}

\begin{solution}{exo:m3:forward-curves-storage-and-convenience-yield:6}
Un jour~: $0.004 \times 50 \times 1 = \$0.20$ par baril. Dix jours~: $0.004 \times 5 \times 5.5 = \$0.11$.
Étaler le roll réduit le coût, mais allonge la \omterm{def:m3:forward-curves-storage-and-convenience-yield:index}{fenêtre de roll} pendant laquelle d'autres peuvent prendre position en amont.
\end{solution}

\begin{solution}{exo:m3:forward-curves-storage-and-convenience-yield:7}
40.9~\%, 38.3~\%, 35.3~\% et 33.6~\%~; le rapport du quatrième au premier est de 0.82. Avec $e^{-\kappa \times
0.25} = 0.82$, $\kappa = 0.79$ par an (une demi-vie d'environ 0.88 an).
\end{solution}

\begin{solution}{exo:m3:forward-curves-storage-and-convenience-yield:8}
Le fonds détient des contrats à terme, non du pétrole~: il gagne la variation du spot, plus le \omterm{def:m3:forward-curves-storage-and-convenience-yield:roll}{rendement de roll}, plus les intérêts. Avec la
courbe en contango plus de la moitié du temps, le roll a coûté 0.71~\% par an de 1985 à 2024, et davantage dans
certaines décennies~; les frais et les coûts de roll propres à l'indice s'ajoutent.
\end{solution}

\begin{solution}{pb:m3:forward-curves-storage-and-convenience-yield:1}
\textbf{1.} 10 millions de barils. \textbf{2.} 0.04~\$. \textbf{3.} $0.04 \times 3 = \$0.12$ par baril.
\textbf{4.} 6~millions d'USD. \textbf{5.} 72~millions d'USD. \textbf{6.} Elle vend le deuxième contrat et achète
le troisième, soit le même spread que celui que l'indice va négocier. \textbf{7.} Pendant la fenêtre, revendre le spread
à l'indice~: elle gagne la hausse du spread, jusqu'à 0.20~\$ par baril sur les positions prises avant le premier jour.
\textbf{8.} Capital, marge, risque que le spread bouge pour d'autres raisons, et concurrence d'autres
front-runners. \textbf{9.} Une prise de position précoce fait bouger le spread avant la fenêtre~: l'indice paie toujours,
mais le gain se partage entre davantage de traders. \textbf{10.} Les traders achètent et vendent aux cours du marché sur
une information publique~; ils prennent la contrepartie d'une transaction connue, ce qui relève d'une fourniture de liquidité par anticipation.
\textbf{11.} 0.20~\$ en un jour, 0.11~\$ sur dix jours. \textbf{12.} Transparence et réplicabilité~: les produits
répliquant l'indice doivent le reproduire exactement. \textbf{13.} Les contrats éloignés se situent là où la courbe est
plus plate, de sorte que le \omterm{def:m3:forward-curves-storage-and-convenience-yield:roll}{rendement de roll} en contango est moins négatif. \textbf{14.} En comparant les rendements des spreads à l'intérieur
et à l'extérieur de la période de roll, de 2000 à mars 2010~: un impact de prix statistiquement significatif, du front-running
ratios de Sharpe pouvant atteindre 4.39, et environ 3.6~\% du rendement annuel perdu par les investisseurs indiciels. \textbf{15.} Un
rendement plus faible que le spot ne le laisse penser, du fait du \omterm{def:m3:forward-curves-storage-and-convenience-yield:roll}{rendement de roll} et du coût de roll. \textbf{16.} Les deux~: la règle rend
l'indice réplicable et peu coûteux à faire tourner, mais fait payer à l'investisseur l'impact. \textbf{17.} En contango, le nouveau
contrat est plus cher que l'ancien~: la position se reporte chaque mois sur des contrats plus onéreux.
\textbf{18.} La forme de la courbe (\omterm{def:m3:forward-curves-storage-and-convenience-yield:roll}{rendement de roll}), la règle de roll de l'indice et son impact, les frais et le
\omterm{def:m3:forward-curves-storage-and-convenience-yield:roll}{rendement du collatéral}. \textbf{19.} \emph{Résultat nommé~:} 0.12~\$ par baril et par roll, 6~millions d'USD par mois et 72
millions d'USD par an dans le modèle illustratif. \textbf{20.} Celui qui fournit le stockage ou l'inventaire que
la courbe rémunère~: les détenteurs de la matière première en backwardation, les stockeurs en contango.
\end{solution}

\begin{solution}{iq:m3:forward-curves-storage-and-convenience-yield:1}
Le bénéfice implicite de la détention de la matière physique~: l'option de l'utiliser immédiatement. Il est élevé quand
les inventaires sont faibles et qu'une rupture de stock est coûteuse, ce qui se traduit par une backwardation.
\iqlookfor{dépendance à l'inventaire et forme de la courbe.}
\end{solution}

\begin{solution}{iq:m3:forward-curves-storage-and-convenience-yield:2}
\omterm{def:m3:forward-curves-storage-and-convenience-yield:roll}{Rendement spot} (variation du prix de proximité), \omterm{def:m3:forward-curves-storage-and-convenience-yield:roll}{rendement de roll} (convergence des contrats détenus, positive en
backwardation) et \omterm{def:m3:forward-curves-storage-and-convenience-yield:roll}{rendement du collatéral} (intérêts sur le cash de marge).
\iqlookfor{trois composantes, et le signe du \omterm{def:m3:forward-curves-storage-and-convenience-yield:roll}{rendement de roll}.}
\end{solution}

\begin{solution}{iq:m3:forward-curves-storage-and-convenience-yield:3}
Les inventaires et l'offre réagissent aux chocs au fil du temps, de sorte que les chocs sur les échéances proches sont plus forts que sur les éloignées~:
l'\omterm{def:m3:forward-curves-storage-and-convenience-yield:samuelson}{effet Samuelson}. Modéliser le spot comme processus à retour à la moyenne (Ornstein--Uhlenbeck) plus un facteur lent~: chaque
contrat à terme a une volatilité égale à $\sigma e^{-\kappa\tau}$ plus celle du facteur de long terme.
\iqlookfor{retour à la moyenne et courbe à deux facteurs.}
\end{solution}

\begin{solution}{iq:m3:forward-curves-storage-and-convenience-yield:4}
\omterm{def:m3:forward-curves-storage-and-convenience-yield:floating}{Stockage flottant}~: acheter le contrat proche, affréter un pétrolier, vendre le contrat éloigné. Risques~: hausse des taux d'affrètement avec la demande,
le financement, les pertes opérationnelles, et le rétrécissement du spread avant de pouvoir le verrouiller.
\iqlookfor{le carry trade et ses coûts.}
\end{solution}

\begin{solution}{iq:m3:forward-curves-storage-and-convenience-yield:5}
Étude d'événement~: rendements des spreads les jours de la fenêtre comparés aux autres jours, en contrôlant le niveau de la courbe et
la saisonnalité, sur de nombreux mois~; tester si l'effet varie avec la taille de l'indice~; fenêtres placebo.
\iqlookfor{une étude d'événement avec contrôles et placebos.}
\end{solution}

\begin{solution}{iq:m3:forward-curves-storage-and-convenience-yield:6}
Date de roll (échéance, bascule de volume, jours fixes), ajustement (aucun, différence, ratio), quel contrat
détenir, et le traitement des prix négatifs. Pièges~: ajustement par ratio sur les prix négatifs, prix ajustés a posteriori
que personne n'a jamais négociés, et des signaux calculés sur des niveaux ajustés.
\iqlookfor{règles de roll, choix d'ajustement et leurs artefacts.}
\end{solution}
